\chapter{万历：整饬之后，战事、财政与辽东（1573—1620）}
\label{ch:wanli}

\section{共享编年叙事：考成、清丈与一条鞭}

\section{报送、边防与河道（1573--1577）}

\eventanchor{MM-1573-KAOCHENG-LAW}{万历元年，张居正推行考成法}，要求部院与地方官按期限申报、完成政务。它把“办到哪里”变成可催问的问题，却没有让北京由此看见每个州县的执行。实录、奏疏和后出的制度材料能确认命令和报送的中枢线；地方文书才可能在特定地点显示期限、工程和征解怎样成为实际责任。

同年，\eventanchor{MM-1573-WANG-GAO-CAMPAIGN}{李成梁进攻王杲集团，王杲被擒后由明廷处置}。辽东卫所、边市和建州诸部彼此牵动，王杲的失败改变了一组首领与关系网，却不能把建州社会预先写成后金的雏形。明方处置意在重整边防与交涉条件；部众如何选择依附、贸易或迁移，仍取决于当地道路和力量。军报可确证行动，不能替诸部划出固定阵营。

\eventanchor{MM-1577-PAN-JIXUN-RIVER-CONTROL}{万历五年潘季驯受命整治黄河、保运漕路}。所谓束水攻沙不是离开地方的技术口号：河道约束到哪里，哪里就要筹料、出役、修堤，并承担决口后风险重新分配。同年\eventanchor{MM-1577-ZHANG-JUZHENG-DUOQING}{张居正父丧期间奉旨夺情留任}，政务连续与礼制要求之间的冲突进入公开争论。河防工程与夺情争议看似无关，却都显示改革的行政速度要靠具体人手、文书和地方服从维持。

\section{田亩、讲学与征解的差异（1578--1582）}

\eventanchor{MM-1578-EMPRESS-WANG-INSTALLATION}{万历六年册立王氏为皇后}，完成幼帝成年后的重要礼仪程序；同年\eventanchor{MM-1578-NATIONWIDE-SURVEY}{朝廷令各省清丈田亩、核查隐漏}。中央命令、各省启动和部分报亩可以相互参照，丈量的完成时间、办法和争议却不相同。一张田册把土地变为可比较、可征收的对象，也把业主、佃作、书手和里甲之间的冲突写进新的行政劳动。

\eventanchor{MM-1579-ACADEMY-RESTRICTIONS}{万历七年限制讲学与书院活动}，将地方士人的声望与结社纳入朝廷的警惕。禁令能说明国家试图约束什么，不能证明书院因此消失；碑记、文集与地方志只能为少数地点补充后续痕迹。\eventanchor{MM-1580-ONE-WHIP-ORDERS}{次年朝廷推动合并田赋、丁役与杂派的征收办法}，\eventanchor{MM-1581-ONE-WHIP-IMPLEMENTATION}{又在中央督促下扩展一条鞭的实施}。银米折纳、市场、借贷和旧有差役在各地的组合不同，不能把政策扩展写成同年完成的全国实测。

\eventanchor{MM-1582-ZHANG-JUZHENG-DEATH}{万历十年张居正去世}。死亡与丧仪可以确认；将其后的每项变化全归因于个人缺位，反而会遮蔽已存在的皇帝、内廷、部院和地方关系。清丈、考成和一条鞭留下的是可催问、可登记的工具，不是已经消除执行差异的结果。

\begin{debate}[清丈的命令不等于各县完成]
《神宗实录》、相关奏疏与《明史·食货志》维持中枢时间线；地方赋役册、黄册残页、契约、河工碑志和地方志只能在局部追问执行。命令、定额、报数、实收与家庭负担不能互相替代，材料覆盖较好的地区也不能代替全国。
\end{debate}

\subsection{期限走到州县以后}

考成法先改变的不是田亩本身，而是事情在衙门里等待的方式。部院的公文有日期，省、府、州、县的回报也要能接回这个日期；一项河工、一笔钱粮、一次勘灾或一件积案，若迟迟不能结案，便不再只是地方官署内部的拖延，而可能成为上级可以追问的失期。北京由此取得的是一条催问的链，而不是一架能越过驿路、河道和乡村、直接观看办事过程的眼睛。县令收到限期，仍须依靠典吏抄录旧案、书手核对册籍、里甲传唤人户；其中任一环的缺漏，都可能在送达上级之前被改写成一份看似齐整的答报。

这也改变了地方官的选择，却没有把选择取消。河决时，官员若先调集堤料和役夫，便要把尚未备齐的文书留待稍后；若先把报表补得合格，河岸的险段未必会等候。征解时，催纳较易的户口可以使账面先有进项，逃亡、讼争或无力折银的人户却可能被留到后面。并非每一次迟延都是欺瞒，也并非每一次如期具报都等于事情已经办妥。考成把这种先后次序置于考核压力下，使官员更须说明何事可办、何事待办、责任由谁承担；它不能替代地方社会对田界、劳力、价格和道路的判断。

\eventref{MM-1573-KAOCHENG-LAW}{万历元年的考成}因而同辽东的军情发生了并不显眼却很实际的联系。李成梁处理王杲集团时，明廷能够接到的是边将、巡按和部院所能写成的军报、请赏与处置意见。辽东城堡之外还有边市、翻译、部众往来和不同首领的利益，它们不必以同一速度进入北京。把王杲被擒写成一个确定节点是必要的；把此后所有建州部众的依附、迁动和贸易都写成该节点的自动后果，则会把文书的终点误当成边地关系的终点。报送制度能使中枢更快追问一份军报，不能让它先于当地人知道哪一处市道、山口或亲属关系会改变下一次交涉。

这层限度在后来尤其重要。辽东不是在万历末年才突然需要粮饷的边线；早年的卫所、市场和首领关系已经要求军政人员在防守、招抚和交换之间分配有限资源。考成能够缩短一份文书在途中的停留，却不能凭空增加可以久驻的兵丁、可供转运的马匹，或愿意按既定价格出粮的商人。万历初年把责任逐级写清，留下的是日后调度可以借用的格式；它也留下一个难题：当边务、河工和钱粮同时具报时，哪一件先被裁决，仍取决于人事、信息和可动用的库藏。

\subsection{河堤、粮船与折纳之前的成本}

万历五年治河时，这个难题落在黄河和运河相接的堤岸上。潘季驯所面对的不是一条可以按图纸安置的水线。水势、淤沙、旧堤、支渠、村落和漕船共同限定工程；堤身要用土、木、石与束草，材料须从附近或更远处调来，役夫则要离开本来的农事、运输和家庭。朝廷的任命与工程主张可以确证一项方案被提出并取得支持，不能使每一段堤防都拥有相同的材料、工时与维护。所谓保运，首先是让一批粮船在可通行的季节抵达下一道闸口；它从来不是河患、决口和沿河负担已经消失的保证。

漕粮北上和地方赋役在这里相互牵连。粮食由产地入仓、过闸、登船、转运到京师，沿途每一次停泊、验收和换载都要由人、船和文书接续。河工占用的劳力与物料，可能同地方本来要承担的差役重叠；水路不畅，又会使本应按期抵达的粮食、银两或军需滞在途中。对京师而言，河防是维系首都供给的工程；对岸边居民而言，它还可能是取土、修堤、迁避和等待水势的日常风险。工程碑记往往保留捐资者和完工的时刻，地方志也可能记下决口与修复；两类记录都不能替代那些未留下姓名的役夫和小户的全程经验。

张居正的夺情争议在这个时候显得格外尖锐，并非因为河工、边防和礼制原本是一件事，而是因为它们都依赖连续的协调。支持留任者担心中枢在清丈、河防和边务相继推进时失去能够汇总部院意见的人；反对者以居丧制度质疑这种把政务连续置于礼制之前的安排。两种担忧都不能被削成一条简单的派系口号。留任并不使河水按计划行走，也不使地方官自动同意征役；居丧之议也不只是脱离事务的道德姿态。它使朝廷必须公开回答：为了维持一个可催问的行政链，可以把多少决定集中在一个人的任期与权威上。

这种集中带来的便利很快会同一条鞭的征解问题相遇。若漕运、河工和边镇仍有大量实物需求，银两的折纳就不能被理解为把一切供给都变成了轻便的货币。银到达州县税柜之后，仍须兑换、采买、交库或转解；银价、米价、脚价和沿途关津的费用，都会决定同一笔名义款项能换得多少粮、布、木料或雇工。某些地方的市镇和商路使这种兑换较易进行，另一些地方却要在实物、差役、借贷与临时摊派之间拼凑。改革要减少名目的纷杂，未必能减少资源由乡里走向河道、仓储和边地时所受的摩擦。

\subsection{清丈册中的人，和不能写进册中的人}

万历六年的清丈把这一摩擦推进到田界。丈量者必须在地面辨认田亩，在册上安排亩数、业主、等则与应负担的项目；田地却可能有典卖、分家、租佃、侵占、荒废和长期未改的旧称。对有契约、宗族关系或地方声望的人家，怎样申报一块田并不只是技术问题；对佃作者、小户、流寓者和已经逃避旧役的人，册上写下或未写下某个名字，也会改变谁被追索。地方书手熟悉旧册和本地称谓，因而是清丈不可或缺的中介；这种中介性同时意味着中央所得到的“可比”数字，已经经过地方知识、利益和协商的整理。

所以，清丈的行政成绩不能只从一串报亩数字判断。某省先报，或许表示丈量与造册较早完成，也可能表示汇总口径较早固定；某县迟报，可能与抗拒、灾荒、地形、书手不足或上级催办的节奏有关。徽州和江南保存下来的契约、册页能够让我们看见一些田土交易、租佃和支付期限，却正因保存条件特殊，不能替代北方军屯、沿河灾区或西南山地。它们最有价值之处，不是替帝国提供一条平均数，而是提醒读者：同一项清丈命令抵达不同地方时，面对的是不相同的产权语言、市场机会和地方权力。

\eventref{MM-1578-NATIONWIDE-SURVEY}{清丈}同年伴随皇后册立，二者都需要礼仪与文书，却把国家能力投向不同的对象。册立以宫廷礼制确认皇帝成年后的家内秩序，田册则试图把分散在村落、户籍和旧账里的土地转化为可征解的条目。前者的完成可在典礼程序中较明确地辨认；后者即使有了中央命令和各地报送，也还要在随后征收中反复接受检验。把这两件事放在同一年，不是要把家庭礼仪化为财政技术，而是说明朝廷能够以不同形式制造秩序，而每一种秩序都有自己的传递距离。

万历七年限制书院讲学，也应放回这条传递距离中理解。士人通过书院、家塾、同年和地方交游讨论政务时，可能批评考成的压力、赋役的安排或官员的作为；中枢担忧的正是这种声望与议论不完全听命于官学和任命系统。禁令能够说明朝廷想收束什么，不能告诉我们每一处书院是否关门，或地方士人是否从此不再组织知识和互助。有人转入家塾、文会或其他名义下的交往，有人则因地方条件、经费和官员态度而沉寂。讲学的限制与清丈的推行看似一紧一松，实则都显出同一个边界：国家可以规定可接受的形式，却不能独自决定地方关系如何继续运作。

\subsection{一条鞭不是一条已经铺好的银路}

到万历八、九年，合并田赋、丁役与杂派的推动给清丈所得的册籍安排了新的用途。用田亩作为较主要的计征基础，折纳一部分银两，可以把原来分散在不同名目、不同服役日期中的责任重新编排；它并不意味着旧有差役、人丁负担、实物供应和地方变通在某一天全部消失。地方官要决定哪些项目可以折银，何时征收，如何解送；书手要把旧额转成新式账目；里甲和人户还得在收成、价格和借贷条件下凑齐可用的支付。制度文本规定的是方向，征收过程却总要在某个县的仓房、集市、渡口与农忙时节落脚。

银两在这条链中不是抽象的改革成果。江南的丝织、棉布和手工业，福建、广东的港口贸易，内地的粮食、盐与山货，都让某些市镇更易取得银；这不等于银会沿着一条无阻的路线流向所有乡村。美洲白银经马尼拉等海上枢纽进入东亚，是许多商人、船主、翻译、港口管理者和季风航路共同造成的流动，途中还有走私、海盗、货价和殖民政权的风险。它可以扩大某些交换的媒介，却不能保证一个需要缴纳折银的家庭恰在征期拥有合适成色和数量的银两。海外贸易的存在更不能证明沿海之外的州县都以同样方式进入市场。

对一个缺少现银的人户而言，折纳的时间比名目本身更直接。丰年或市价有利时，出售粮食、布匹或手工业品也许能换得银两；歉收、道路受阻或价格不利时，同一笔税额可能迫使家庭借贷、预卖产品，或请求把责任延后。对地方官而言，提前解出部分银款能够回应上级催办，却可能把最困难的欠额留给后来的追征；对商人与有储蓄者而言，代纳、放贷和转运又可能带来新的位置。我们不能由少数契约断言所有家庭如何承受，也不能把地方间的差异写成改革无效。能够确认的是，\eventref{MM-1580-ONE-WHIP-ORDERS}{万历八年的推动}和\eventref{MM-1581-ONE-WHIP-IMPLEMENTATION}{次年的扩展}把征收责任更紧地接到土地、市场和中介之上，而这三者从未在全国同样分布。

从河堤回看，银路也不是脱离运输的。军粮、漕粮和河工所需的实物仍要走水陆道路；银若要购买这些物资，也要经过商人、仓储和地方市场。海上贸易增加的金属货币并没有替代运河的闸坝、辽东的马匹或乡村的劳力。正因为如此，后来朝鲜战事需要的补给不能只从“国家已有一条鞭”推出。制度提供了更便于核算和催解的若干办法，能否把这些办法转成跨海、渡江和前线所需的粮饷，仍取决于不同地方是否能够筹得、换得、运得并按时交出资源。

\subsection{首辅去世以后，工具仍在，时间已经改变}

\eventref{MM-1582-ZHANG-JUZHENG-DEATH}{张居正去世}并没有使田册、限期或折纳的做法从各地账簿上立即抹去。已经形成的征收办法会在州县、府库和地方社会中继续被使用、修改或抵触；改动其节奏的，是皇帝、内廷、部院、督抚与地方官之间不再有同一位首辅持续撮合的条件。此后对张居正名位和家产的处置可以显示朝廷政治评价的转向，却不能代替逐处核验清丈或一条鞭怎样存续。把改革的延续全部归功于旧制，或把其后的困难全部归咎于首辅已死，都会跳过制度真正要经过的人、物和道路。

这种时间差在万历二十年以后变得格外沉重。宁夏、朝鲜和其后的西南、辽东并不是同一种战争，所需的路线、粮源和地方关系也不相同。朝鲜方面的援军要面对港口、海峡、渡江、朝鲜地方的粮食与避难人口；明廷可以下令调兵和题拨，不能因此使船只、粮食与银两按同一速度抵达。中朝日的记录可以互校某些出兵、会战和撤退的先后，却分别从本国朝廷、战场和海路观察战争。兵力、斩获和民众遭遇并不因后来归入一个“援朝”名称而变得一致。对明朝财政而言，战争的后果不止于一笔可以在京师账上结清的支出，而在于地方仓储、运输者、军户与商人被反复要求接续缺口。

朝鲜战事之后，海上网络也没有只留下胜败的回声。福建、广东与海外港口的商路继续影响银和商品的可得性，却同时暴露于季风、海防、港口规制与殖民暴力。万历三十一年马尼拉华人社群遭到镇压的事件，只能说明一处特定港口和社群所遭遇的危险；它不允许把所有华人航海者写成同一种群体，也不能由不同档案给出的数字重造精确总数。它却提醒我们，所谓全球贸易并非一条自动向内地输送白银的管道。港口的分类、船只是否出航、商人能否返航，都会影响某一地区在何时能够把货物换成银，而这又会在征税和借贷的时刻传到地方社会。

到万历后期，宫廷的继承争论和辽东的军情使同一套调度能力承受更紧的挤压。慈庆宫的案件和围绕继承的奏疏首先是中央程序，不能从审讯记录推断每个官员或京师居民的意图；但任命、批答和责任归属一再迟延，会改变边报、饷项和地方请求得到处理的先后。努尔哈赤建国、辽东攻势与熊廷弼经略防务，把城堡、道路、粮食、马匹和情报网络重新推到朝廷面前。明方、满文和朝鲜材料各有自己的政治语言，不能据单一战报画出无争议的控制线；仍可看见的是，边地所需的资源必须穿过已经承担河运、地方征解与其他军务的行政通道。

\eventref{ML-1619-901}{辽饷的增筹}正是在这个背景下进入各省。会计书、部院奏疏与实录能让我们分辨预算、题请和制度安排，却不能把它们直接换算成每一户实际缴纳的数额。某地先解的款项未必已到前线，账上的欠额也未必能说明哪一个家庭首先失去粮食或土地。真正连续的链条，是从万历初年限期报送、清丈造册和折纳征收开始，经过河道、市场与海港的不同速度，最后在辽东城堡和地方税柜之间被拉得更紧。它没有预先决定明朝的结局，却解释了为何万历末年每一次新的边警、运输阻塞或宫廷迟延，都会使已经相互牵连的财政与社会压力更难分开。

\subsection{从资格、讲席到税柜：一条不走直线的地方政治链}

\eventref{ML-1584-901}{张居正身后官爵被追夺、家产被籍没}，先在北京的诏令、处分和人事记忆中出现；它随即也给仕途上正在等待任用、考选或外放的人留下一个很实际的问题：一套曾经用来催问政务的做法，今后还会由谁解释、谁来承担与旧首辅相连的名声。不能把这次处置写成全国官员同时改变了行事方式。实录和后出的列传所能确知的是朝廷如何处理其名位，难以逐一记录一名候选官员、一位府学教官或一个州县书手由此作出的选择。然而，地方治理并不只由已经在任者维持。通过学校、考试、荐举、同年与幕府进入仕途的人，要学习怎样读公文、怎样写呈文，也要从前辈的褒贬中判断何种勤敏会被视为能干，何种依附会在政局转向时成为危险。政治评价因此不是京城中一段与州县无关的余波；它会改变地方行政人员安排自己的声名、请示与人际凭借的尺度。

这种影响不宜被说成某种整齐的“士大夫意见”。府学和县学有不同的经费、学田、教官与生员来源，乡试、会试留下的名次和籍贯也只使一部分读书人进入后世可见的名册。名册能确认谁在某次考试中取得资格，却不能告诉我们一个地方有多少家庭读书、多少塾师讲授，或一个新中者是否立刻拥有支配本地事务的力量。许多准备应试的人仍在宗族、田产、店铺、书坊和亲友资助中筹措纸笔、路费与居停；有些人终身未能进入可被中央记录的官员行列。正因如此，科举不应被想成一架从乡村直接升到朝堂的梯子。它是把经典语言、书写技能和可被官府承认的资格带入不同地方的渠道之一，而地方上能否供养这条渠道，又取决于土地收益、捐助、学校田租和城市市场。

\eventref{ML-1587-901}{海瑞卒于南京任所}使这种资格与地方声望的关系获得一个可见的节点。南京官署中的去世、墓祠和后来的传记并不是同一种材料：前者把他置于任官与礼制程序，墓祠和地方记忆则会选择何种德行值得被保留。后人将海瑞塑为清官，并不能倒过来证明南京官场或江南州县普遍按同一准则运作；同样，也不能因纪念带有褒扬而否认它会成为地方读书人讨论官箴、家产和为政边界的资源。一处祠祀、一篇祭文或一部方志所显出的，是某些捐资者、撰文者和地方官愿意公开承认的海瑞，不是一个没有分歧的社会共识。可是，它们说明仕途声望可以在官署之外继续流通，并在学校、宗族和城市的读书圈里为后来的行政行为提供可援引、也可争议的尺度。

这与万历初年的考成压力相接，却不是同一件事。考成所要的是期限、责任与报送；应试和讲学所训练、检验或传播的，则包括把地方经验译成经义、策论、公文和交游语言的能力。一个县令处理丈量、催科或河工时，未必亲自从书院取得方案；但他所倚重的幕友、书手、学生和地方绅士，常常在学校、文会、宗族或旧日同年关系中彼此可辨。这样的关系能够帮助搜集旧例、筹办劝捐、解释一项征解，也可能让某些人的意见较早进入县衙、另一些人的困境较难被写成可上报的文字。我们很难由存世文集追完一件案件如何穿过每一层关系，尤其难以看见不识字者如何评价这些安排；因此不能把“士绅参与”写成天然的地方自治。能够把握的是，国家要求更细密的册籍和期限时，地方社会中掌握文字、信用和往来渠道的人，往往成为把命令化为具体程序的中介。

\eventref{ML-1596-901}{矿监税使的派遣}让这种中介位置骤然敏感。矿区、关津与市集原已有地方官、书吏、牙行、铺户、矿工和行旅所熟悉的称量、保结、运输与纠纷处置；内廷使者进入以后，何物可被指认为应税，谁有资格验看货物，哪一笔银应先解出，便不再只由常规的州县链条决定。地方官若据旧例申明困难，可能是在维护本地生计，也可能是在争取自身的行政空间；商人投诉关卡阻滞，可能揭示实际损失，也可能服务于请求豁免或竞争对手间的争执。奏疏把这些冲突带到朝廷时，常以弹劾、请撤或请恤的形式排列事实。它们很适合让人看见某种请求为何被提出，却不能直接告诉我们每一座市镇征到多少、每一名矿工受了何种影响。

不过，税使的出现仍改变了地方政治的时间。常规田赋可以依农时、册籍和历年催科安排先后，商税与矿税的验认却常在货物正在流转、商人正待过关的时刻发生。对于一位依赖市集出售布匹、纸张、茶货或山货的人，等待验看和补办保结会同价格、脚费和债期一起决定一趟交易是否还有余利；对于县衙而言，若急于满足上级或内廷使者的期限，原来用于治安、赈济、河工或诉讼的人员和文书便可能被挪去应付新的征取。这里不能把一次派遣解释为一条从矿场直达国库的银路，更不能把地方的每一项冲突都归因于税使。它所显示的是：当征取权绕开或重排既有层级时，熟悉地方交易规则的人既可能成为协助者，也可能成为申诉、规避和抗衡的组织者。

这也是考试资格与商业社会相遇的地方。取得生员、举人或进士资格并不使人天然站在市场之外；许多读书人与家族的田租、借贷、典当、商号、书坊或捐资有不同程度的联系，而地方富户和商人也可能资助学校、桥梁、寺观和刻书，以取得声望、信用或与官府交涉的语言。具体联系必须由契约、碑刻、账簿或地方志逐处核查，不能由一个人的功名或一座书院推定其财富来源。特别是保存较多的江南、徽州材料，只能显示那些地点的产权和文书习惯，不能代替北方军镇、西南山地或沿海渔村。不过，当赋役更常需要折银、军需又扩大对转运和信用的需求时，读书、捐助和经商所构成的关系网便不再是财政之外的装饰：它们影响谁能临时垫付、谁能写成一份申诉、谁有熟人可以在府城或省城递送文书。

\eventref{ML-1601-901}{利玛窦获准居北京}，使另一种以文字、礼物、翻译和接见为条件的流通进入这条链。居京首先是一项特定的准入，不等于一种宗教已经在全国获得同样的空间。传教士书信、士人文集和后出的记述各自保留了交往的片段，也各自有向不同读者安排人物位置的理由。我们可以看见利玛窦与中国士人、官员接触的某些踪迹，却不能从这些精英材料推断普通居民如何理解天主教，更不能把皈依、好奇、礼物交换和对器物或历法的兴趣合为同一个态度。对北京的官员和读书人而言，陌生知识能否进入谈论，取决于引荐者、语言、书写方式和宫廷许可；对地方读者而言，一部书或一段消息能否抵达，又取决于刻印、携带、赠阅和可负担的纸墨。

\eventref{ML-1602-901}{《坤舆万国全图》的刊行}恰好让这种差别具体化。地图上出现的海洋、国家和地名，并不自动改变一名读书人处理本县赋役或一名商人判断一季货价的办法；刊刻是由译述、书写、绘制、校对、纸张、雕版和出资共同完成的工作，读图则还要经过书坊、馈赠、藏书、讲席与转抄。存世版本和题跋能帮助辨认某些刊刻与拥有者，却尤其容易保存有能力署名、题写和收藏的人。它们不能量出有多少人见过此图，也不能证明观看者接受了同一种世界秩序。地图的政治作用不在于立刻替代地方经验，而在于它让一部分人可以用新的空间词汇谈论海路、朝贡、传教、贸易与边界；这些词汇若进入策论、书信或官场交游，才可能在行政和知识之间发生更慢、更不均匀的往返。

宗教传播也应从这种往返来理解。寺观、祠祀、家礼和地方善会原已构成不同的仪式、募缘和互助空间；天主教知识的进入并没有把它们一举清空，也不能只被写成“中西冲突”的预告。某位士人参与讨论、某处有人受洗、某部文本被重印，分别是可以追查的不同事实；它们不能证明全城信仰转换，更不能告诉我们没有留下文字的人如何在祭祀、疾病、丧葬、借贷和迁徙中理解宗教。反过来说，若只因材料多出自精英便把这些交流视为无关紧要，也会看不见知识如何通过翻译、书写和结交改变可被提出的问题。地方社会所面对的不是抽象的“西学”或“传统”，而是具体地决定一笔捐资投向学校、寺观还是刻书，决定一名读书人通过哪一种师友关系取得新书与新名声的选择。

\eventref{ML-1603-901}{马尼拉华人社群遭西班牙殖民当局镇压和屠杀}使知识与货物的流通显出它们共有的脆弱性。马尼拉不是一座只为中国白银服务的港口；殖民当局、华人商贩与工匠、船员、翻译、教会人员和来自不同海域的买主，都以不同身份被港口制度分类。西班牙档案可能保存统治者如何报告秩序与暴力，福建地方材料或后起记述可能保留另一端得到的消息；译名、人数、死者身份和消息传到中国的路径并不完全一致。故而不能用一个常见数字替所有受害者计数，也不能把他们概括为没有职业、籍贯与亲属差异的单一“华人”。可以较谨慎地说，港口暴力会切断或改变船、货、信用和返乡消息的既有安排，而这些安排原本已与福建沿海及更远内地的债务、订货和家计相连。

这种影响也不等于某一次海外惨案立刻改变了帝国税额。货物是否装船、商人能否归还本钱、银两在何处结算，先由季风、航路、殖民规制和合伙关系决定；它们随后才可能以价格、汇兑、借贷和纳税能力的变化抵达市镇。一个等待货款的人家也许先感到的是债期无法兑现，一名经营纸墨或布匹的商人也许先感到买主缩减，地方税柜却未必立刻留下相应的记录。我们没有覆盖所有港口、所有行商和所有家庭的连续账簿，不能把海上风险直接换算成某省的财政损失。可是，正因为折银和军需都要求市场能够在特定时日提供货币、货物与运输，海外联系的中断不应被放在地方财政叙事之外。它告诉人们，商业信用既能让资源跨越很远的距离，也能把远方权力的暴力带入一纸借约、一笔欠款和一次不能按期起解的交易。

\eventref{ML-1604-901}{顾宪成等重修东林书院并讲学}之后，书院可以更清楚地被看作资格、地方声望和公共议论相互折合的场所，而不应被预先涂成一个全国一致的党派总部。重修需要场地、经费、主持者与地方认可；讲学需要听众、文本、师友和来往道路；声望则要通过文集、碑记、书信、同年或门人的传播才可能超出无锡一地。碑记和文集最能显示主持者如何说明其目的，也最容易抹去未被邀请、不同意或无力参与的人。即使在书院附近，商人捐资、宗族支持、地方官态度和生员之间也未必形成同一种政治立场。把所有这些人立即归作“东林”，会比当时的地方关系更整齐。

但书院的政治意义也不只在后来被贴上的标签。读书人讨论吏治、赋役、官箴和名节时，可以把一处县政、一件征收或一项任命放进较广的评价语言；这些讨论若借助刻书、传抄和同年网络流动，便可能影响谁被信任、谁的奏疏值得转递、哪一项地方请求得到声援。它们并不直接替州县作决定。县令仍须面对田界、仓储、逃亡、讼案和催解，商人仍须估量货价与路险，佃户和小民也不因讲席上的议论而拥有同等发言机会。书院能提供的是一种联结：它使地方经验有机会被写成道德和制度问题，也使中央的人事与财政争论能够借由地方名望找到新的听众。联结越强，围绕它的争执也越可能影响官员对声名、交游和奏报风险的判断。

到\eventref{ML-1613-901}{楚宗案进入处置程序}时，这种判断再次被宫廷的宗法与法律问题牵动。宗室谱系、科道言论和审理程序首先属于朝廷所能看见的范围；它们不提供一幅社会对继承问题意见一致的图像。可是，案件既要经过奏疏、批答和官员的责任分配，便会使熟悉章奏、律例与名分语言的人更关切何种陈述能够进入程序。地方学校、书院和同年交游未必直接参与个案，不能被虚构为一张操纵京师的网；但它们是这些语言得以学习、争论、抄传和评价的若干场所。宫廷程序若久而未决，也会使任官、人事与请示的节奏更难预期，而这种不确定并不会停在宫门以内。

\eventref{ML-1615-901}{梃击案的审讯}把这层不确定集中到慈庆宫的门禁、侍卫、供词与审理。史料能较为明确地追踪案件怎样被报告、谁被讯问以及朝廷如何处置；它却不能从一份供词推出所有参与者的动机，也不能让后来的派别叙事替代当时人的犹疑。对京师官员而言，继承安全同任命、批答和责任归属相连；对外地的官员和读书人而言，京师所发出的文书节奏会影响一份请示何时得到回应、一项人事是否有人裁决。这里所谓影响不是说每一个州县问题都由梃击案造成，而是说行政链本来就要依赖可预期的接收、批答和承当者，宫廷程序的紧张会使这条链更容易积压。

随后，\eventref{ML-1616-901}{努尔哈赤建国号金}、\eventref{ML-1617-901}{熊廷弼经略辽东}和\eventref{ML-1618-901}{辽东攻势}把积压已久的问题重新压到奏报、任官和地方征解上。边报需要能够读写、核验、转递并承担责任的官员与书吏，城堡需要能按时到达的粮、马、火器和修缮材料；然而，一张功名名册不能保证某地立即有合用的人手，一处书院的声望也不能替城堡运来粮食。考试、讲学和刻书所形成的知识网络，只能提供部分人才、语言和信用，不能消除道路、库藏、市场和战场信息的限制。正是在这层限度中，地方政治的文化面与财政面才真正相接：谁能把边情写成被重视的文书，谁能为临时筹解找到保人、商路或捐资者，谁又因旧债、地域和身份而留在可见性之外，都关系到一项防务要求如何被改写成地方上的实际负担。

万历四十七年朝廷增筹辽饷，遂把这条不走直线的链收回到明确的编年节点。朝廷所能列出的，是筹饷、分派和催解的项目；州县所必须处置的，却是现银从哪里来、谁愿垫付、道路何时可走、旧欠和新派如何并列。科举资格和书院声望不能把这些困难化解为道德问题，海外贸易和商业信用也不能保证银两恰在征期出现；宗教与新知识的交流同样不会自动改变一座税柜的存银。它们仍是万历地方社会的重要条件，因为文字、交游、刻印、信誉与航路决定了信息、货物和请求可能穿过哪些人。到了辽东军情加急的岁月，国家既需要这些分散的能力，又无法把它们完全纳入同一张账册。由考成后的地方资格、矿税下的市镇交涉、北京与无锡的知识流通再回到辽饷，可以看见的不是一条必然导致危机的直线，而是一套越来越依赖中介、信用和有限时日的行政世界。

\subsection{税柜、山路与边堡}

\eventref{ML-1589-901}{明廷授予努尔哈赤龙虎将军等衔}时，辽东已不是把京师命令直线投向边外的空地。这一名号是朝廷能够确认的一次封授，却只能说明明廷试图用官衔和既有交涉语言安置一位首领，不能证明山口、村寨和市集中的关系就此稳定下来。辽阳、抚顺及其间的城堡需要驻军、修缮与日常给养，城外的互市、翻译和往来人户又使消息、皮货、布帛与粮食在不同的规则下移动。守堡官员希望把道路视为可调度的线，行旅、商人和边民却首先要衡量一段路能否通行、货物能否验放、债主是否等待结算。一个官衔可以把某位行动者写进朝廷秩序，不能替代这些每天决定交换能否继续的条件。后来边事急转时，朝廷面对的正是这种早已存在的差异：它拥有可以发给的名义，却不拥有每一处关口的信用、熟路的人手和可随时取用的物资。

万历二十四年的\eventref{ML-1596-901}{矿监税使派遣}把财政的催取推入另一类道路。矿山、关津和府城并非从此才有货物流转；地方官、书手、牙行、脚夫、铺户和小贩原已在称量、保结、寄存和转卖中形成各自可依循的办法。税使到来后，哪些货物要验看、由谁出具保结、何时可以放行，往往同一笔买卖的时限纠缠在一起。对正等待把纸、布、茶货或山货送往市镇的人来说，耽搁并不只是账面上一项税名：价格可能在候验时变化，雇来的脚夫和牲口仍要付钱，先前借入的本钱也有到期的日子。对州县来说，接待、稽察和申报会占用本来用于仓储、讼案、赈济或治安的文书与人手。于是，朝廷希望取得的银两，并不是从一处矿口自动流向京库，而要经过许多能否相互承认的印信、货价和人际担保。

这种情形不能被写成“商人”与“国家”各站一边的整齐对峙。经营长途交易的人或许有机会借熟悉道路、货源和汇兑而代为垫付；受雇转运的人却可能只在货物延误时承担损失；地方富户既可能在劝捐、借贷和保结中获得与官府交涉的机会，也可能因新的验认被卷入原来不必处理的风险。各地奏疏保存的多是请求裁夺、弹劾侵扰或申明困难的声音，最容易让后人看见冲突如何被送往北京，却未必留下那些没有能力上呈的人怎样改换生计。地方志、契约和商账能够把少数地点的路程、价款或债务写得更近，保存较多的地区又不能替代所有山场和集镇。因此，矿税的区域效果只能随着具体地点辨认：有的地方首先感到的是货物被留在关口，有的地方是征取权与县衙旧例相撞，不能把一份奏报推成全国市场同时停顿的证明。

财政上的这种紧张，与西南山地的军务在同一时期相互逼近，却不应被简单说成一项政策必然引出另一场战争。\eventref{ML-1599-901}{杨应龙在播州再起兵}以后，播州的山路、寨堡和邻近州县把军事问题变成了川、黔之间具体的通行问题。山地用兵所需的不只是兵册上的人数，还包括能够翻越险阻的向导、负载粮米和器械的人畜、临时停驻的地点，以及在雨水、山道和地方关系变化中仍能接续的消息。朝廷在\eventref{ML-1600-901}{任命李化龙筹划播州用兵}时可以部署督抚和军队，却不能由任命本身决定每一支队伍从何处得到粮食、哪一处关隘可以按期通过，或沿途居民将以何种代价提供夫役、草料和居处。对于附近乡村和市镇，军队经过既可能带来采购、雇佣与转卖的机会，也可能带来征发、涨价、逃避和旧债难清；两种结果在不同道路和不同家庭中可以同时存在。

\eventref{ML-1606-901}{海龙屯攻破、杨应龙死而播州改流}之后，前线的攻守并没有把这片山地变成一张已经安静的行政图。改流意味着新的官署、文书和征收原则要在旧有聚落、道路和权力关系中寻找落脚处；它所要求的，不只是把城寨纳入地图，更是让命令能抵达、纠纷能被受理、税粮或差役能被说明。军队撤离后的运输者、手工业者、山货经营者和被征调的住民，未必以同一方式看待新的设置。明廷叙事容易把攻城与改置排成胜利的终点，地方材料与遗址所能保存的则往往是更零碎的居住、道路和遗存线索。它们不能替我们补写每一户的遭遇，却足以阻止把“平定”当作资源、人口和市场立刻恢复原状的同义词。西南的这一经验提醒人们，军事行动的区域后果不只在战报里，也在后来谁有资格调用劳力、谁能在新设衙门前陈述地界与债务、谁必须绕开改变后的关隘和税卡。

到\eventref{ML-1601-902}{朝廷下令撤回部分矿监税使}时，征取链条同样不能被想成一纸撤令便全部复原。撤回的是一部分使者和某些明确的处置，并不等于此前已经发生的验货、拖欠、保结与官民争执会同时消散。一个商号是否能够重新取得货源，要看合伙人是否仍有本钱、路上的债务是否结清、地方官是否愿意恢复旧例；一户因征发或滞销而借贷的人家，也未必能从北京的诏令中立刻得到现银。反过来，也不能把撤罢写成毫无作用：中央对征取方式的调整会改变地方官、商人和中介者预期谁有权验看、谁应承担责任。政策的来回正显示财政能力的限度不在于有没有一条命令，而在于命令要经由多少人、地点和时间才会成为可辨认的日常秩序。

这种限度在辽东重新显得尖锐。万历四十四年的\eventref{ML-1616-901}{努尔哈赤建国号金}不是此前受封关系的机械反面，也不能把不同部众、城堡居民和贸易参与者预先合成两个永远对立的阵营；它却使明廷必须以新的紧迫感处理边外政权的组织与辽东防务的空缺。次年\eventref{ML-1617-901}{熊廷弼受命经略辽东}，经略所面对的不是一张等待填色的边图，而是城堡守备、仓储、马匹、修缮材料、翻译消息和地方人心彼此牵连的局面。若一座城的粮秣不能按时接续，守军的选择会受制于饥乏和欠饷；若道路上的车辆、驮畜和人力先被别处征用，仓中有名目的物资也未必能在需要时变成前线可用的东西。辽东境内的卫所家庭、屯田住户、城镇铺户和往来客商不是战场之外的背景，他们分别承受着守城、供给、避险和交换秩序改变的后果。

\eventref{ML-1618-901}{后金向明辽东发动攻势}以后，这些约束不再能只靠延后答复来容纳。城堡失守或告急会改变附近居民的去留与市集的安全，来往受阻也会使原本依赖城镇的粮食、布匹、牲畜和信息难以按原来的方式汇集。明方军报、满文记录和相邻政权的观察各自把行动者置于不同政治语言中；它们可以相互校正某些先后，却不能据一方的胜败叙述断定每一处道路、市场和人群已经怎样归属。尤其不应把“辽东”当成一块没有内部差别的边区：抚顺一带的城防压力、辽阳周围的仓储与交易、远离前线的州县收到的征解要求，经过的都是不同的距离和中介。前线需要的是及时的粮、马、器具和能够守住岗位的人，后方感到的却可能是货价、夫役、借贷和催解先后发生的改变；两端由同一项军务牵连，并不表示两端经历了同样的危机。

因此，\eventref{ML-1619-901}{辽饷的增筹}不宜只被读作京师向各省加上一项数字。县衙要把新旧责任放入既有的征收日程，府与省要面对解送、核销和欠额的排列，能够提供现银或组织陆路转运的人则要估计信用、货价、道路和回款的风险。这里的商人网络并非一支随时听令的后备军：有人可能因熟悉货源、保人和路线而成为筹措的一环，有人也会因资金被占、货物受阻或债期逼近而无力继续。对乡村而言，承担未必总以同一形式出现；土地收益、手工业收入、雇佣机会和亲族借贷不同，接受催解的能力也不同。会计材料可以让我们辨认朝廷如何设想收支，部院文书能够保留谁在何时请求何种办法，却不能把预算、题请和实际到达的物资压缩成同一条事实，更不能替每个家庭算出一份共同负担。

征解真正改变的，常常是地方上原先分开的几种时间被迫重叠。农户要等收成、出售和结算，手工业者要等原料、订单与买主，商号要等外地伙伴还款，县衙却须按上级所定日期将银两凑成可以起解的数目。若市场上的现银一时紧张，先有信用的人可能以借贷、代纳或押货维持周转；没有可供抵押的田产、店铺或亲族关系的人，则更可能在同一征期里先卖掉粮食、预支劳作，或让旧欠叠到新欠上。这里不能把任何一种反应断言为普遍经验。存世契约通常留下能够订立和保存文字的人，地方账簿也往往只记录已经进入衙门程序的款项。然而它们提示了一个不能省去的中间层：赋役并不是从一张官方名册直接落到辽东前线，它要在家庭收支、坊市买卖、保人担保和州县文书之间一次次改变形态。

区域的差异使这一中间层尤其重要。播州山地的难处在于险路、聚落分散和军政改置尚未完全稳固；辽东的难处则在于城堡、边情和军事需求不断改变物资应当送往何处、由谁接收。两地都需要人力、粮食、牲畜、工具和可靠消息，但不能把西南用兵后的征发方式移作东北守备的缩影。离前线较远的府县也不会因此置身事外：它们可能首先遇到的是催解文书、商路上货价与信用的波动，或地方官为筹措而重新寻找捐助、保结和暂借的办法。反过来，辽东城镇的需要并不只在边墙以内消耗，它会经由采购、雇佣、欠饷和亲属迁移，改变周围居民对市场、卫所和官府的判断。所谓“后方”并不是没有战争的空间，而是承担后果的方式与城堡不同的空间。

这也解释了为什么不能只用白银是否流入、某省是否富庶来判断辽饷能否兑现。银两能够作为税额和支付手段，并不保证它恰在征期、恰在可通行的地点、恰在愿意承担风险的人手中出现。一个经营货物的人也许有账面资产，却因客商未归或货物未售而拿不出即刻可用的现银；一个州县即使先行起解，也仍要面对途中保管、交收凭据和后续补足的难题。相反，短期筹措成功也不能证明基层负担已经被公平分摊，或边堡所需必然按时到达。把这些不确定写清，既不是把晚明财政说成完全失灵，也不是把市场想成能够自动修补国家缺口的机器。它所呈现的是一套必须不断借助地方知识、关系与风险承担者才能运行的征收体系，而边事越急，留给这些人调节的时间便越少。

在这一过程中，官员的选择也不能只按“加派”或“反加派”两种姿态分类。有人会优先使一笔款项尽快离开本县，以免上级追责；有人可能要求缓期、分次或另寻来源，以免把压力集中在某一时点；有人借助本地富户和商人取得短期周转，也有人因此让掌握信用者获得更大的发言权。每种做法都可能缓和眼前的一处缺口，同时把风险移到另一个群体、另一个月份或下一任官员手中。奏疏中的理由未必等同于真实动机，地方传说中的贪廉评判也未必能说明账目的全貌；但从一份请求为何要写成那样、由谁为它作保、它避开或触及了哪一类资源，仍可看出财政命令在地方被重新安排的痕迹。万历末年所显露的并非单纯的中央命令减少或市场力量增加，而是两者都更依赖那些能够把资源、文字和时间接合起来的中介。

从辽东的封授与边市，到矿税下的关津和山场，再到播州山路上的动员，以及最后回到各省税柜与辽东城堡之间的催解，可以看见万历后期并非只有一条由中央向边地伸出的财政线。国家依赖地方官的文书、商人的信用、转运者的劳力、城镇的仓储和居民在压力下作出的选择；这些条件又受季节、道路、旧债、地方权力和战争消息不断改写。它们共同解释了为何新的边警会如此迅速地扩大为区域性的难题，却不授权我们把结局倒写进每一笔征收和每一次交易之中。
\begin{figure}[htbp]
  \centering
  \begin{tikzpicture}[
    x=.88cm,y=.88cm,font=\small,
    hub/.style={draw=timeline!85,fill=paper,rounded corners=2pt,align=center,inner sep=3pt},
    court/.style={hub,draw=dynasty!90,fill=dynasty!14},
    region/.style={hub,draw=source!80,fill=source!12},
    note/.style={font=\scriptsize,align=center,text=source}
  ]
    \fill[paper] (-.35,-.4) rectangle (12.15,6.35);
    \fill[dynasty!10] (.15,.05) rectangle (11.9,5.95);
    \node[court] (beijing) at (6.15,5.1) {北京\\户部、仓储与军政拨给};
    \node[region] (tongzhou) at (7.8,4.1) {通州\\仓储、转运节点};
    \node[region] (xuzhou) at (5.25,2.8) {徐州、淮北\\河工与漕运风险};
    \node[region] (jiangnan) at (3.0,.95) {江南\\田赋、漕粮与市镇网络};
    \node[region] (huizhou) at (6.85,.85) {徽州及内陆\\田亩登记、契约与折银};
    \node[region] (liaodong) at (10.2,3.1) {辽东边镇\\军饷与补给需求};
    \node[region] (northwest) at (1.35,3.65) {西北、宣大等边镇\\军需与转输压力};

    \draw[->,ultra thick,color=timeline!90] (jiangnan) -- node[left,note]
      {漕粮、船户与河道维护} (xuzhou);
    \draw[->,ultra thick,color=timeline!90] (xuzhou) -- node[left,note]
      {运河、闸坝与转运（示意）} (beijing);
    \draw[->,very thick,color=dynasty] (beijing) -- node[above,sloped,note]
      {军费、命令与文书调度} (liaodong);
    \draw[->,thick,color=dynasty] (beijing) -- node[above,sloped,note]
      {支给与核销的行政链} (northwest);
    \draw[->,thick,dashed,color=source!90] (huizhou) -- node[below,sloped,note]
      {清丈、赋役折银与地方中介} (jiangnan);
    \draw[<->,thick,dashed,color=source!90] (huizhou) -- node[below,sloped,note]
      {税额、报解与实收不能互推} (beijing);

    \node[draw=source!75,fill=paper,rounded corners=2pt,align=center,
      font=\scriptsize,text=source] at (9.55,.72)
      {箭头显示资源与文书的关联方向，\\不表示实际运量、税额或通航状态。};
  \end{tikzpicture}
  \caption{万历时期漕运、河工与财政征解的关联路线（示意）。潘季驯治河、清丈与一条鞭式赋役整顿，可据《神宗实录》、河渠志、《明史》〈食货志〉及《河防一览》追踪其政策与工程脉络；黄册、鱼鳞图册、徽州契约、地方志和河工档案用于校验不同地点的登记、征收与运输。江南至北京的线路和边镇节点仅展示相互依赖，不能据此推定某年漕粮实到、折银比例、军饷数额或所有地区的实际负担。}
  \label{fig:vol05:transport-fiscal-routes}
\end{figure}


\section{改革的余波与三条战线（1584--1598）}

\eventanchor{ML-1584-901}{万历十二年追夺张居正身后官爵、籍没其家产}。朝廷处置的是一位已故首辅的名位与家产，不能替代对考成、清丈和一条鞭在各地如何延续的核验。\eventanchor{ML-1587-901}{万历十五年海瑞卒于南京任所}，南京地方志、墓祠和列传留下的是任官与纪念的不同线索，也不应被写成改革的简单终局。

\eventanchor{ML-1589-901}{万历十七年明廷授努尔哈赤龙虎将军等衔}。这一封授属于辽东边务的行政关系，尚不能倒写为后来战争不可避免的开端。三年后，\eventanchor{ML-1592-901}{哱拜及其部众据宁夏叛乱，明军平定}；同年\eventanchor{ML-1592-902}{丰臣政权入侵朝鲜，明廷开始援朝调兵}。宁夏城、黄河水利和西北粮道是一条战线，朝鲜的港口、渡江、军粮和地方动员是另一条；京师必须把兵员、道路、饷项和不同方向的军报接进同一套决策。

\eventanchor{ML-1593-901}{万历二十一年李如松等明军与朝鲜军收复平壤}。平壤不是孤立捷报：援军登陆、渡江、补给与朝鲜地方动员共同构成战区。\eventanchor{ML-1596-901}{万历二十四年矿监税使被派往各地征办矿税、商税}，矿区、市场和州县遂成为内廷派员、地方官、商民和奏报相遇的地点；派遣令不能证明各地征解相同。\eventanchor{ML-1597-901}{日军再度侵朝后，明廷续遣援军，由邢玠经略军务}；\eventanchor{ML-1598-901}{丰臣秀吉死后日军撤离，露梁海战发生在撤退阶段}。中、朝、日材料共同维持先后线，却在兵力、战果、译名和民众遭遇上保留不同观察位置。

\section{播州、矿税与海上网络（1599--1606）}

\eventanchor{ML-1599-901}{万历二十七年杨应龙在播州再起兵}。海龙屯周围的山路、关隘与寨堡比北京命令更先决定军队能到哪里；四川、贵州及邻近地区必须让兵员、粮饷和转运接得起来。\eventanchor{ML-1600-901}{朝廷任李化龙总督筹划播州用兵}，显示中央选择将有限财力投向一条山地战线。明方战报的军数、伤亡和“改流”目标有请饷、问责和胜利叙事的压力，不能写成山地社会已同意何种治理。

\eventanchor{ML-1601-901}{同年利玛窦获准居北京}，标志一段可以核查的特定交往；\eventanchor{ML-1601-902}{朝廷又下令撤回部分矿监税使}，只能证明撤罢程序与政策转向，不能将先前负担折成全国平均线。矿税奏疏、地方志与食货材料须区别题请、派遣、报数和可见反应。\eventanchor{ML-1602-901}{次年利玛窦与李之藻刊行《坤舆万国全图》}，居京、刊图与读图的范围是不同问题；耶稣会书信、存世版本题跋和李之藻文集能够确认交往与刊刻，不能量出社会传播。

\eventanchor{ML-1603-901}{万历三十一年马尼拉华人社群遭西班牙殖民当局镇压和屠杀}，提醒人们海外港口并不只是知识交流的背景。西班牙档案、福建地方材料和后出的中国记录能互校事件，不能把一个港口的档案分类或死亡数字扩大为所有华人海上活动。\eventanchor{ML-1604-901}{顾宪成等重修东林书院并讲学}，书院碑记、文集和地方志所见的是士人组织及自我表述，不是已被证明的全国政治集团。\eventanchor{ML-1606-901}{万历三十四年明军攻破海龙屯，杨应龙死，播州改流}；设施、道路和地方记忆可由遗址报告和贵州方志辅助定位，实际控制范围仍不能照录单一战报。

\section{继承、辽东与地方征解（1613--1620）}

\eventanchor{ML-1613-901}{万历四十一年楚宗案中的宗室与科道争论进入处置程序}。实录、\textit{国榷}和奏疏能显示案件怎样进入制度入口，不能仅凭供词裁定策划者动机。\eventanchor{ML-1615-901}{万历四十三年张差持梃闯入慈庆宫，引发梃击案审讯}。宫门、侍卫和审讯文书使继承焦虑成为具体行政程序；案卷能说明谁被问讯，不能替所有士人、宦官和京师居民安排态度。

\eventanchor{ML-1616-901}{万历四十四年努尔哈赤建国号金、建元天命}；\eventanchor{ML-1617-901}{次年熊廷弼受命经略辽东}。到\eventanchor{ML-1618-901}{万历四十六年所谓“七大恨”发布、向明辽东发动攻势}，城堡、道路和粮饷从边务问题转为朝廷持续调度的中心。明、满文与朝鲜材料各自陈述这个转折，不能以一方战报固定军力或控制线。

\eventanchor{ML-1619-901}{万历四十七年朝廷为辽东战事增筹军饷，辽饷征解进入地方行政链}。会计书和部院奏疏记录预算、题请与制度语言，不能直接化为每户实纳或全国平均负担。辽东防务、播州之后的西南整合、东南海上联系和江南、北方的征解至此由同一套有限调度能力相互挤压。

\subsection{制度得失：可征解的国家，难以同时调度的边地}
\textbf{所得}：考成和清丈留下的报送与征解工具，使朝廷能够把援朝、播州、矿税和辽东送入可催问的程序。

\textbf{代价}：文书与预算的可见性没有消除道路、地方中介、军饷和边地信息的差异；每条战线都在争夺同一套有限能力。

\textbf{承接}：辽饷的征解与辽东防御，将在天启、崇祯两朝继续挤压地方财政和军事选择。
\section{把“整饬”送到县里：1573--1582}

万历初年，考成、清丈和一条鞭常被放在一条从张居正出发的改革线上；读者若跟着诏令走到州县，看到的却是一系列需要反复接续的工作。\eventref{MM-1573-KAOCHENG-LAW}{考成法}要求部院与地方按期限报送和完成事务，使迟延成为可被追问的行政问题。它提高了中枢看见某些工程、赋役和官员行为的能力，也使地方官更有理由以符合上级格式的语言书写本地情形。报告越整齐，并不必然越接近村庄、河堤或军屯的全部现实；它首先显示哪些事情能进入催办的文书链。

辽东与河道说明这种文书链并非只服务于财政。\eventref{MM-1573-WANG-GAO-CAMPAIGN}{对王杲集团的处置}发生在卫所、边市与建州诸部相互牵动的地方。首领被擒能够改变一组关系，不能将建州社会预先写成后金国家的必然起点。\eventref{MM-1577-PAN-JIXUN-RIVER-CONTROL}{潘季驯整治黄河、保运漕路}则让国家能力落在堤防、闸坝、役夫和沿河家庭身上。束水攻沙是一种工程方案，不是一句能保证河流听命的口号；河道在哪一段被约束，哪一段就要筹料、出役、承担决口和淤塞风险。辽东的军情与运河的保运由不同材料留下，却共同争夺粮、银和人手。

\eventref{MM-1577-ZHANG-JUZHENG-DUOQING}{张居正夺情留任}使政务连续与礼制要求直接相撞。父丧期间是否应离职，关乎一位首辅的个人伦理，也关乎考成、清丈和边防协调是否失去能够汇总文书的人。争论并不说明所有官员只站在两种固定立场上；有人忧虑礼制、有人忧虑行政中断，也有人借此重新安排人事关系。\eventref{MM-1578-EMPRESS-WANG-INSTALLATION}{册立皇后}和\eventref{MM-1578-NATIONWIDE-SURVEY}{全国清丈}同在万历六年，前者以宫廷礼仪完成皇帝成年后的家庭秩序，后者试图把田亩、隐漏和赋额写进更可比较的册籍。两项事务都提醒我们，国家以不同方式把人和物纳入可见性。

清丈的困难不在于中央能否发出命令，而在于田界、佃作、租佃、隐占、灾荒、旧账和地方权威怎样被写进或排除在登记之外。业主希望保护既有利益，佃作者和小户可能面对不同的差役转移，书手和里甲要在上级口径与本地关系之间寻找可行写法。存世田册、契约与州县材料能让少数地点的争执显形，不能给出全帝国的平均故事。\eventref{MM-1579-ACADEMY-RESTRICTIONS}{限制书院讲学}也应放在这种环境里理解：诏令能说明中央担心地方士人如何结社和议论，不能证明知识网络因此消失，或所有书院都以同样方式顺从。

\eventref{MM-1580-ONE-WHIP-ORDERS}{合并赋役的推动}和\eventref{MM-1581-ONE-WHIP-IMPLEMENTATION}{一条鞭在多地扩展}，把丁役、田赋和杂派更紧地接到银两与地方征解上。银不是均匀流入每一个乡村的财富标记：银米比价、运费、借贷、商路、实物缴纳与旧有差役都会改变一项折纳对家庭意味着什么。某地能够以银纳税，未必表示其市场充足；某地保留实物或劳役，亦不等于不在改革压力之内。制度文本、预算、报数、实收和家庭负担必须分开，才不会把改革写成一张立即覆盖全国的账单。

\eventref{MM-1582-ZHANG-JUZHENG-DEATH}{张居正去世}后，报送和征解工具仍在，能够协调它们的人事与政治条件却改变了。皇帝、内廷、部院、地方督抚和州县各自重新估量责任，任何后来的松动、追夺或地方差异都不能简单归为首辅不在。改革留下的真正难题是：国家已经学会要求更细密的说明，却仍要面对说明能否变成粮、银、河工和地方服从的问题。

\section{战争怎样改变财政的方向：1584--1598}

\eventref{ML-1584-901}{张居正身后被追夺官爵、籍没家产}，使政治评价和财产处置进入同一程序。朝廷可以重写一位首辅的名位，不能由此抹去清丈、考成和一条鞭已经进入各省的文书与地方实践。\eventref{ML-1587-901}{海瑞卒于南京}也不应被处理为另一种改革纪念碑：南京官署、地方士人、墓祠和后出的传记各自保存一部分形象，无法证明某种清廉政治已经遍及江南。两件事都显示，制度的后果常通过人事、财产和地方记忆被重新解释。

\eventref{ML-1589-901}{授努尔哈赤龙虎将军等衔}时，辽东的关系仍以卫所、贡市、首领、翻译和边官交涉为媒介。封授说明明廷希望在既有行政框架内安排一位地方行动者，不能倒写为后来的战争已经注定。\eventref{ML-1592-901}{宁夏叛乱}与\eventref{ML-1592-902}{丰臣入侵朝鲜、明廷援兵出动}却让不同空间同时要求军费。宁夏围城关系到城墙、黄河、粮道和西北军镇；朝鲜战场关系到海运、渡江、朝鲜地方动员和跨国情报。北京需要在同一套预算与命令中处理两者，地方承担者却感到的是不同道路上的征发和不确定。

\eventref{ML-1593-901}{平壤收复}说明援朝军可以在特定条件下与朝鲜力量共同改变战场，不能把它写成战争的结束。粮食、病疫、港口、朝鲜民众、日军补给与明军轮换持续决定行动。中、朝、日材料各自记录的部队名称、地名、日期和战果并不完全一致；差异不是可删去的噪声，而是不同政权如何看待战场和民众的证据。\eventref{ML-1597-901}{再度援朝}到\eventref{ML-1598-901}{日军撤离、露梁海战}，使战争的收束同丰臣政权内部变化、海上撤退和地方破坏交缠。不能只用“援朝成功”覆盖朝鲜社会所承受的迁移、征粮和重建。

同一时期的\eventref{ML-1596-901}{矿监税使派遣}把财政从部院常规渠道伸向矿区、集市和州县。内廷派员、地方官、商人、矿工和居民围绕税额、开采、运输与执法相遇；一纸派遣令不说明各地征解相同，地方奏疏也可能强调最有利于请撤、请恤或问责的一面。\eventref{ML-1601-902}{撤回部分矿监}是政策转向的程序，不是对先前负担的自动清算。不同地区的矿税、商税、交通和市场条件需要保留自己的尺度，不能折成一户全国平均负担。

\section{西南、海港与知识的不同全球线：1599--1606}

\eventref{ML-1599-901}{杨应龙在播州再起}以后，海龙屯周围的山路、寨堡和跨省转运决定明军能够怎样行动。\eventref{ML-1600-901}{李化龙筹划用兵}不是将一支大军放到空白地图上，而是要使四川、贵州与邻近地区的粮、兵、向导和地方中介接续。\eventref{ML-1606-901}{攻破海龙屯、播州改流}能够确认军事与行政结果，不能说明山地社会已经以同一方式进入新秩序。遗址考古、地方志和实录各自说明设施、记忆与朝廷处置，军数、伤亡和控制范围仍须保持限度。

\eventref{ML-1601-901}{利玛窦居北京}和\eventref{ML-1602-901}{《坤舆万国全图》的刊行}让另一类跨区域联系进入士人和宫廷的文本世界。居京是一项准入，刊图是一项由李之藻等参与的书写、刻印与交换，读图则取决于购买、赠阅、书院和私人网络；三者不能混成“西学传遍中国”。耶稣会书信、题跋和存世版本有助于确认特定交往，也带着作者为欧洲和中国读者安排自身角色的目的。

\eventref{ML-1603-901}{马尼拉华人社群遭屠杀}则提示全球海路同时是殖民暴力、港口分类和生计风险的空间。西班牙档案、福建材料和后出的中文记录能在部分时间与地点上互校，死亡数字和群体组成却不能从一类档案扩大。\eventref{ML-1604-901}{东林书院重修讲学}同样不是一个天然全国化的政治集团诞生；碑记、文集、刻书和地方声望说明一些士人怎样组织讨论和互相承认，不能把所有地方读书人、官员或商人装进同一党争标签。海外港口、北京地图和江南书院都依赖纸张、航路、金钱与中介，构成万历世界中彼此相关却不能相互替代的流通。

\section{继承与辽饷：1607--1620}

万历后期的迟延不能只用“怠政”两个字解释。奏疏、任命、矿税、军饷和继承程序仍在产生，谁能促成决定、谁愿承担后果、哪一份边报能进入宫门，却越来越没有可预期的节奏。\eventref{ML-1613-901}{楚宗案}使宗室、科道和法律程序的争执公开化；\eventref{ML-1615-901}{梃击案}又让东宫安全、供词和继承焦虑集中在慈庆宫门外。实录、国榷和奏疏能说明案件进入何种程序，不能由一份口供裁定策划、动机或所有人的态度。

\eventref{ML-1616-901}{努尔哈赤建国号金}、\eventref{ML-1617-901}{熊廷弼经略辽东}与\eventref{ML-1618-901}{对明辽东发动攻势}，将边政、军事和财政接成更紧的链。辽东城堡需要兵员、火器、粮食、马匹、道路和能够传递情报的地方网络；明、满文和朝鲜材料各自用不同的政治语言叙述行动，不能固定一条无争议控制线。\eventref{ML-1619-901}{辽饷增筹}使这种防务压力穿过江南、北方和各省的征解体系。预算、定额、实解、拖欠和家庭负担不是同一件事；会计书可以记录某部门的口径，不能使每一个县的实际支付透明。

到万历末年，朝廷并非毫无工具：考成遗下的报送习惯、清丈和赋役整顿提供了询问地方的格式，运河、仓储、卫所和市场仍能使资源移动。问题在于这些工具必须同时服务辽东、地方灾荒、官员人事、港口贸易与继承危机。行政能力不是拥有多少条命令，而是能否让不同命令在有限人手、有限信息和不同地方社会中不互相毁损。下一阶段的失城、三饷与流动作战，正是在这种被长期拉紧的条件上展开。
\section{战争经费怎样进入地方生活}

万历援朝最难的部分不在于把军队送到朝鲜，而在于让远离战场的财政与运输长期保持。辽东、山东、直隶与江南的仓储、船只、骡马、布匹和银两都被编入战争的时间表。\eventref{ML-1592-901}{壬辰援朝}的朝廷记录能够列出调兵、饷银和战报，不能把“饷”视作已经抵达士兵手中的实物。银的成色、道路的阻塞、层层经手和地方先行垫付，都会改变一笔军费的实际价值。战争越久，中央越依赖地方官、商人、军户与运输者在命令之外补足缺口，补足又可能以借贷、拖欠和加派的形式留在地方。

朝鲜战场也不能只是明军的舞台。朝鲜王朝、地方义兵、农民、僧侣、日军与明军在不同地区拥有不同信息和目标；首都、港口、山道与粮仓的控制随战局改变。中朝史料各自强调本国政治与行动者，后来的民族叙事又会选择性放大某些贡献。把材料并置，可以确认联盟中的协作、误解和资源争夺，却无法让我们为每个村庄重建完整经验。可以确定的是，战争摧毁耕作与居住的能力，迫使人口迁移，也使东亚海陆交通的风险重新分配。

平壤、碧蹄馆、蔚山等战事的命名容易制造一串清楚的节点，补给的节奏却很少如此清楚。冬季、山地和长距离运输限制行军；军中疾病、逃亡、欠饷和指挥权争论同样限制胜利能否被扩大。军事报告常用斩获与城池呈现成果，失败和困难则可能被推给天气、友军或下属。阅读时应将战果放在材料立场中，不把单一军报转写为无条件的战场全貌。援朝的意义在于它把明朝既有的财政、边防和海上问题同时推入极限，而非只在海外留下光荣或挫败的纪念。

\section{银、税与地方裁量}

战争结束后，中央并未回到一张未受损的账本。矿监、税使与加派是在财政缺口、内廷需求和地方商业增长之间出现的手段。\eventref{ML-1596-901}{矿税使派遣}在制度上扩大了直接取财与监察的渠道，在地方则常通过关卡、市场、矿场、寺观和富户的纠纷显形。官书和奏疏留下大量控诉，也带有反对者的政治位置；不应据此假定所有地区受害同样，或每一项征收都被完整执行。可以肯定的是，税使制度改变了谁能够决定何为应税、谁能申诉、谁须先行缴纳的权力关系。

白银流通使这种关系更复杂。江南的丝织、棉布和手工业，福建、广东的海上贸易，内地的粮食与盐业，都把地方经济接入更广的兑换网络。马尼拉贸易和美洲白银进入东亚的路径，不能被简化成一条直线：中介商人、港口管理、走私、季风、海盗和价格变化共同决定银如何到达某个县的税柜。\eventref{MM-1581-ONE-WHIP-IMPLEMENTATION}{一条鞭法推行}确实提高了赋役折银的比重，却没有取消实物、劳役和地方差异。对有市场渠道的人，折银可能带来弹性；对现金不足或在歉收中被迫变卖的人，它也可能加重风险。

地方官在这套变化中并非只有服从或反抗两种选择。他们可以拖延、折衷、借助士绅筹款、调整征解次序，或把难以承担的项目转入基层。地方档案和方志偶尔保留这些动作，但大多只在争议激化时显现。所谓“地方执行”不是一个中性的末端，而是政策真正成为家庭负担、商机或冲突的地方。中央以考成和诏令要求整齐，实际运行则依赖无数没有写进中央史书的协商。

\section{晚万历的治理时间}

国本争论的长期化，改变的不只是不立太子的礼制问题。继承不确定使官员任用、派系结交和奏疏攻防都带有对未来朝局的计算。册立太子虽给出一个正式节点，却没有立即消除前期积累的猜疑。不能把东林、浙党或其他群体当作固定不变的现代政党；它们是在选官、言官、财政、人事和道德语言中形成又分化的关系网络。派系标签能帮助识别争论，不能替代逐案考察。

播州用兵与辽东紧张说明帝国在不同边缘同时被要求投入。\eventref{ML-1600-901}{播州平乱}需要穿越山地的军粮、土司关系和地方路线，战后又涉及设治、移民和资源重新分配。辽东的女真诸部联合则把互市、封授、边堡和军费问题重新拉紧。\eventref{ML-1619-001}{萨尔浒之战}前后的信息并不完整，后见之明不能被投射为当时人人已经知道的结局。万历末年最明显的不是某项制度突然消失，而是决策、拨款、任命与回应危机的时间越来越难以同步；这种不同步为下一时期的财政、军事和政治失衡留下了条件。
\section{信息、灾害与等待回应的人}

万历后期的行政困境还表现在消息的不同速度。辽东的军情、贵州的战报、江南的税案和福建的海船消息，都要经过地方衙门、驿路、部院与内廷才能成为可裁决的奏疏。路程、抄写、衙门积案和人事关系决定何时被看见；得到批答之后，命令还要沿相反方向回去。面对突发灾害或边警，地方官往往不能等待完整指示，只能在仓储、乡绅、军队和市场之间先作安排。后来文本把某次处置写成忠或奸、勤或惰，实际选择常受到账目、道路和信息不全的约束。

灾荒尤其显示这种约束。旱涝、蝗灾和疫病在不同地区以不同节奏发生，赈济诏书只能说明中央承认危机并下达原则。开仓需要库存和运输，减免需要核实田亩，流民安置需要土地、粮食和当地接受能力。地方志记下灾异与善举时，常以恢复秩序为终点；逃亡、借贷、卖产与家庭离散则更难持续记录。不能因而把救荒写成徒有其名，也不能把一纸诏令当作已经抵达每一位饥民。

晚万历的文化和宗教交流也在这种不均衡中发展。传教士、士大夫、书商和工匠围绕历法、数学、地理与礼仪发生接触，留下的文本多来自能够书写和保存论争的人。它们能够显示知识如何被翻译、选择和争论，不能代表整个社会已经接受或拒绝某种新知。全球白银、海上商品和新知识进入中国，与乡村赋役、城市手工业和宫廷财政同属一张并不平整的网络。到1620年前后，这张网络仍在产生财富与机会，也在每一次银价、战费、灾害和任命延迟时暴露其脆弱。